\section{Additional Details and Results of Experiments}
\label{sec:exp_details}
\ifdefined\msom\else See \Cref{sec:ci} for how point estimates and confidence intervals are calculated for the estimators in \Cref{sec:experiments}. 
\fi

\subsection{Experiments for \citet{KangSc07}}
\ifdefined\pnas
\subsubsection{Data Generating Process}
\label{sec:ks_dgp}
The \textit{true} outcome and treatment  mechanisms to depend on covariates $\xi \sim N(0, I) \in \mathbb R^4$ and $\varepsilon\sim N(0,1)$
\begin{align*}
Y & = 210 + 27.4 \xi_{1} + 13.7 \xi_{2} + 13.7 \xi_{3} + 13.7 \xi_{4} + \varepsilon, \\
\pi(\xi) & = \frac{\exp(-\xi_{1} + 0.5 \xi_{2} - 0.25 \xi_{3} - 0.1 \xi_{4})}{1+\exp(-\xi_{1} + 0.5 \xi_{2} - 0.25 \xi_{3} - 0.1 \xi_{4})}.
\end{align*}
Here, $P[Y(1)]=P[Y\mid A = 1] = 200$ and $P[Y] = 210$ so that a naive average of the treated units is biased by -10. For a random sample of 100 data points, the true propensity score can be as low as 1\% and as high as 95\%. Instead of observing $\xi$, the modeler observes
$$
X_{1} = \exp(\xi_{1}/2),\;X_{2} = \xi_{2}/(1+\exp(\xi_{1})) + 10,\;X_{3} = (\xi_{1}\xi_{3}/25 + 0.6)^3,\;X_{4} = (\xi_{2} + \xi_{4} + 20)^2.
$$
\else\fi
\subsubsection{Implementation Details}
\label{sec:ks_experiment_details}

\paragraph{Linear Outcome Models.} Here, $P_{\rm train}=P_{\rm eval}$, as consistent with the original paper by \citet{KangSc07}\footnote{Although sample splitting could be better, we have chosen to replicate \cite{KangSc07} as closely as possible.}. There is no $P_{\rm val}$ as there are no hyperparameters to tune. 

For the direct method and initial outcome model for AIPW, self-normalized AIPW, and TMLE, we simply regress the dependent variable $Y$ on the observed covariates $X$ using the samples with labels. We use all the samples available to fit a logistic regression model for the treatment variable $A$ (outcome was observed) using the covariates $X$. Following \cite{KangSc07}, in both models, we omit the intercept term when fitting the propensity scores and initial outcome model. Our main insights do not change if an intercept is added to the outcome and propensity models.

We run datasets with 200 and 1000 observations, with and without a truncation threshold of $5\%$. For each configuration, we generate 1000 seeds. 


Table~\ref{tb:linear_simulation_full} displays the results for the bias, mean absolute error, root mean squared error (RMSE), and median absolute error over 1000 simulations with sample sizes equal to 200 and 1000. We also include the results when truncating to 5\% is performed, meaning that $\what \pi(X)$ is truncated to have a lower bound of 0.05. Note that truncation greatly improves the performance of propensity-based methods such as AIPW and TMLE. Without truncation, C-Learner is the best method across all samples when compared to other unbiased methods and is outperformed only by the direct method. When truncation is introduced, the C-Learner is similar in performance to other asymptotically optimal methods.

For the coverage computation, we construct the confidence intervals as described in \cref{sec:ci} and empirically check for each method if it covers the true population mean.
 Table~\ref{tb:coverage_ks} presents the coverage results for a 95\% confidence level. It is worth noticing that coverage results greatly deteriorate when the sample size increases for all asymptotically optimal methods, most likely because both the outcome model and the propensity model are mispecified, and consistency is ensured when at least one of the models is consistent (see \cref{sec:theory_text}).

 Finally, Table~\ref{tb:flipped_ks} displays the results for bias and mean absolute error (MAE) for estimating $P[Y(0)]$ rather than $P[Y(1)]$. In this case, all the asymptotic optimal methods perform very similarly, and all are better than the direct method with OLS. To conclude, we emphasize that the C-Learner is the only asymptotically optimal method that demonstrated good performance across all configurations of data size, truncation procedure, and data-generating process by either achieving comparable performance to other asymptotically optimal methods, or by improving performance, sometimes by an order of magnitude.




\paragraph{Linear Models with Logistic Link.} In order to instanciate the TMLE and the C-Learner using the logistic link, first we normalize the observed outcome $Y$. In particular, for the observed dataset we compute $Y_{\max} = (1+\alpha)\max_{i:A_i = 1} \{Y_{i}\}$ and $Y_{\min} = (1-\alpha)\min_{i:A_i = 1} \{Y_{i}\}$, where the scaling parameter $\alpha$ accounts for the fact that for the observed dataset, both the empirical maximum and empirical minimum are most likely biased. We chose this procedure to match previous work in the literature, such as the one proposed in \citet{VanDerLaanRo11}, and it does not mean to provide unbiased estimators for such estimates. Observed outcomes are then rescaled by computing $\tilde Y_i = (Y_i - Y_{\min})/(Y_{\max}-Y_{\min})$. In the experiments reported in \Cref{sec:experiments}, we use $\alpha = .1$ as this was the parameter used in \citet{VanDerLaanRo11} when analyzing the same dataset. Different values of $\alpha$ ranging from $0.01$ to $0.2$ leads to the same insights we provide in the main text, with the C-Learner-L dominating the MAE obtained with TMLE-L. The default $\alpha$ for the \texttt{tmle} package available in R~\citep{gruber2012tmle} is 0.01. 

\paragraph{TMLE-L.} Abusing notation, define $\sigma(X) := 1/(1+\exp(-X))$, which is the logistic link function. As an starting point, one solve for the fractional logistic regression using the scaled outcomes $\tilde Y$:

\[
\min_{\theta} \left\{ - \sum_{i=1}^{n} A_i\left[ \tilde Y_i \log \sigma(\theta^\top x_i) + (1 - \tilde Y_i) \log (1 - \sigma(\theta^\top X_i)) \right] \right\},
\]
which leads to $\what{\mu}(X_i) = 1/(1+\exp(-X^\top \theta))$. The second step is to define a parametric submodel and find the optimal fluctuation parameter $\epsilon$. Define
$$
\what{\mu}(X_i,H_i,\epsilon) := \frac{1}{1+e^{-\log \tfrac{\what{\mu}(X_i)}{1-\what{\mu}(X_i)}+\epsilon H_i}},
$$
and note that $\what{\mu}(X_i,H_i,\epsilon)$ is a shift of $\what{\mu}(X_i)$ under the logit transformation in the direction of $H_i$. Therefore, under the logistic link, the optimal fluctuation can be fitted by solving
\begin{equation}
\label{eqn:targeting_logistic}
    \epsilon^\star:=\argmin_{\epsilon \in \R} 
    ~~-\frac{1}{n}\sum_{i=1}^n  A_i \left( 
    \tilde Y_i \log \what{\mu}(X_i,H_i,\epsilon) - (1-\tilde Y_i)\log (1-\what{\mu}(X_i,H_i,\epsilon))\right).
\end{equation}

It is straightforward to verify that the first-order conditions for $\epsilon$ imply that the empirical evaluation of \eqref{eqn:ate-pathwise} at $\what{\mu}(A_i,X_i,H_i,\epsilon^\star)$ is zero. Next, the TMLE procedure is executed as usual by evaluating the fluctuated model on the whole sample. To implement TMLE with the logistic link, we use the \texttt{tmle} package available in R~\citep{gruber2012tmle}. 

\paragraph{C-Learner-L.} For the C-Learner-L implementation we solve
\begin{align*}
\theta^\star = \argmin_{\theta} \; - &\sum_{i=1}^{n} A_i\left[ \tilde Y_i \log \sigma(\theta^\top X_i) + (1 - \tilde Y_i) \log (1 - \sigma(\theta^\top X_i)) \right]\;\text{ such that }\\ 
&\sum_{i=1}^{n} \tfrac{A_i}{\hat \pi(X_i)}\left( \tilde Y_i -\sigma(\theta^\top X_i) \right)=0. 
\end{align*}

Note that this is a convex optimization problem with nonlinear constraints. In order to solve it, we use a Sequential Least Squares Programming algorithm available in the R package \texttt{nloptr} \cite{nloptr2023}.


\ifdefined\response\newpage\newpage\clearpage\pagebreak\else\fi
\label{sec:dual_details}
\paragraph{C-Learner-Dual And Covariate Balancing Methods}

We describe the implementation of the C-Learner-Dual for linear propensity models under the logistic link. We also discuss the implementation of the benchmark methods. For all models, we also implement their self-normalized version, in which the propensities sum to one.


\paragraph{CBPS} We implement the CBPS approach of \cite{imai2014covariate} using the \texttt{CBPS} package available in R~\citep{gruber2012tmle}. We use a linear specification and the over-identified model that is optimized for both propensity score fitting and the balancing constraint with the two-step procedure, following the same procedure as in the package documentation for the KS dataset.


\paragraph{Linear Fluctuation of Propensity Scores} We implement a one‐step calibration of the initial propensity scores via a low‐dimensional fluctuation. The procedure resembles in spirit the parametric submodel approach of TMLE for outcome models for achieving the semiparametric variance bound. We follow \cite{tan2010bounded}'s construction for $\what\mu_{LIK2}$: Define the calibration covariates and form the extended propensity model
    \[
      \omega_i(\lambda)
      = \hat\pi_i + H_i^\top\lambda,
      \quad
      \ell(\lambda)
      = \sum_{i=1}^n\bigl[A_i\log\omega_i + (1-A_i)\log(1-\omega_i)\bigr].
    \]
where $H_i=(1-\what\pi(X_i))\what\mu(X_i)$ and $\pi_i$ is the vector of propensities fitted by standard MLE. We maximize likelihood $\ell(\lambda)$ with respect to $\lambda$ to obtain $\hat\lambda$. In particular, we use the Nelder-Mead method available in the package \texttt{optim}  in R. 

Next, we fix $\omega$ and proceed to the second step, and we solve 
\[\max_{\lambda_{step}}  \sum_i \left[ T \frac{\log(\pi_i + \lambda_{step}' [1,\mu(X_i)]^\top) - \log(\omega_i)}{1-\pi_{i}} - \lambda_{step}' [1,\mu(X_i)]^\top \right].\]
Define $\omega_{step} = \pi_i + \lambda_{step}' [1,\mu(X_i)]^\top$. The first order condition implies that \[
\frac{1}{n}
\sum_{i=1}^n
\Bigl(1 - \frac{A_i}{\omega_{step}(X_i)}\Bigr)
\,\hat\mu(X_i)
\;=\;0,
\]
which is the C-Learner-Dual constraint. 

Similar to the C-Learner, the resulting estimator attains the semiparametric variance bound, and is doubly robust. Our implementation for linear fluctuations differs from \cite{tan2010bounded}'s construction of $\what\mu_{LIK2}$ only in that we omit the $h_2$ component in the first step, which, in their paper, makes the resulting estimator sample-bounded (in the sense that estimator values are within the range of outcome values seen in the sample); this omission is for a more fair comparison between estimators.



\paragraph{C-Learner-Dual} First, we parametrize the propensity‐score model by
\[
\pi_{\beta_{\text{C-Learner-Dual}}}(X_i)
\;=\;
\frac{\exp\bigl(X_i^\top \beta_{\text{C-Learner-Dual}}\bigr)}
     {1 + \exp\bigl(X_i^\top \beta_{\text{C-Learner-Dual}}\bigr)},
\]
and fit \(\beta_{\text{C-Learner-Dual}}\) by maximizing the Bernoulli log‐likelihood on the treatment indicators \(A_i\).  Equivalently, we minimize the negative log‐likelihood
\[
-\ell\bigl(\beta_{\text{C-Learner-Dual}}\bigr)
=
- \sum_{i=1}^n \Bigl[A_i\,\eta_i 
- \log\bigl(1 + e^{\eta_i}\bigr)\Bigr],
\qquad
\eta_i \;=\; X_i^\top \beta_{\text{C-Learner-Dual}}.
\]
Because this objective is smooth and convex, we also supply its gradient in closed form:
\[
\nabla_{\beta_{\text{C-Learner-Dual}}}
\bigl[-\ell(\beta_{\text{C-Learner-Dual}})\bigr]
=
-\,X^\top\bigl(A - p(\beta_{\text{C-Learner-Dual}})\bigr),
\quad
p(\beta_{\text{C-Learner-Dual}})_i
=
\frac{e^{\eta_i}}{1 + e^{\eta_i}}.
\]

Next, to ensure semiparametric efficiency of the resulting inverse–propensity‐weighted estimator, we impose the finite‐sample balance constraint
\[
\frac{1}{n}
\sum_{i=1}^n
\Bigl(1 - \frac{A_i}{\pi_{\beta_{\text{C-Learner-Dual}}}(X_i)}\Bigr)
\,\hat\mu(X_i)
\;=\;0,
\]
where \(\hat\mu(X_i)\) is a fixed outcome‐model prediction.  Rearranging shows this is equivalent (up to an additive constant) to
\[
\sum_{i=1}^n
A_i\,\hat\mu(X_i)\,e^{-X_i^\top \beta_{\text{C-Learner-Dual}}}
\;-\;
\sum_{i=1}^n
(1 - A_i)\,\hat\mu(X_i)
\;=\;
0.
\]
We likewise provide the Jacobian of this constraint:
\[
\nabla_{\beta_{\text{C-Learner-Dual}}}
\bigl[\text{constraint}\bigr]
=
-\,\sum_{i=1}^n
A_i\,\hat\mu(X_i)\,e^{-X_i^\top \beta_{\text{C-Learner-Dual}}}\,X_i.
\]

Because we now have a convex, differentiable objective with one smooth equality constraint, we solve it using an augmented‐Lagrangian scheme. In particular, we use the \texttt{nloptr} package (\cite{nloptr2023}) with the NLOPT\_LD\_AUGLAG solver. Tolerance and maximum number of iterations are set to $10^{-8}$ and 1000, respectively. Finally, the solution $\widehat\beta_{\text{C-Learner-Dual}}$ defines the dual C-Learner propensity estimate
\[
\widehat\pi^C(X_i)
=
\frac{\exp\bigl(X_i^\top \widehat\beta_{\text{C-Learner-Dual}}\bigr)}
     {1 + \exp\bigl(X_i^\top \widehat\beta_{\text{C-Learner-Dual}}\bigr)},
\]
and the average treatment effect is then estimated by
\label{sec:dual_details-end}
$$
\widehat\psi
=
\frac{1}{n}
\sum_{i=1}^n
\frac{A_i}{\widehat\pi^C(X_i)} \;Y_i.
$$


\ifdefined\response\newpage\newpage\clearpage\pagebreak\else\fi

\ifdefined\response\newpage\newpage\clearpage\pagebreak\else\fi

\paragraph{Gradient Boosted Regression Tree Outcome Models.} We demonstrate the flexibility and performance of the C-Learner in which we instantiate C-Learner using gradient boosted regression trees using the XGBoost package \citep{xgb} with a custom objective, as outlined in \cref{sec:methods_boosting}. $\what\pi$ is fit as a logistic regression on covariates $X$ using L1 regularization (LASSO). 

For each seed, and sample size, we randomly take half of the data for $P_{\rm train}$ 
and half of the data for $P_{\rm eval}$. For the first phase of Algorithm~\ref{alg:boosting}, we perform hyperparameter tuning using $P_{\rm val} = P_{\rm eval}$. Hyperparameter tuning is performed using a grid search for the following parameters: learning rate (0.01, 0.05, 0.1, 0.2), feature subsample by tree (0.5, 0.8, 1), and max tree depth (3, 4, 5). We also set the maximum number of weak learners to 2 thousand, and we perform early stopping using MSE loss on $P_{\rm val}$ for 20 rounds. Hyperparameter tuning is performed separately for C-Learner and the initial outcome model.

For the second phase of Algorithm~\ref{alg:boosting}, we use the set of hyperparameters found in the first stage. The weak learners are fitted using $P_{\rm eval}$ and $P_{\rm val} = P_{\rm eval}$. In order to avoid overfitting in the targeting step, we use a subsampling of 50\% and early stopping after 20 rounds.

In Table~\ref{tb:linear_simulation_xgb_trim} we provided additional details with truncation in order to understand the impact in the results of AIPW and TMLE that produced unreliable estimates for sample sizes of 200 and 1000 as well as the impact of truncation in the other estimates. When ``small'' truncation is performed (0.1\%), C-Learner is still better than any other estimator. When truncation to 5\% is performed, TMLE achieves the best pointwise performance, and C-Leaner and AIPW-SN are statically equal in terms of MAE.

Finally, in Table~\ref{tb:coverage_ks_xgb} we present the coverage results for the estimators computed without truncation and a sample size of 200 and 1000 presented in the main text on Table ~\ref{tb:linear_simulation_xgb}. For a sample size of 200, asymptotically optimal methods have similar coverage and substantially lower than the target of 95\%. When the sample size increases, the coverage results of C-Learner deteriorate, although it still presents the best performance in terms of MAE. 


\subsubsection{Estimator Performance When Varying Overlap Between Treatment and Control}
\label{sec:ks_sensitivity}
In order to test the sensitivity of different asymptotic optimal methods with respect to the propensity score, we modify the probability of treatment in the data generation by introducing a scaling parameter $c$ so that
\begin{align}
\pi(\xi) = \frac{\exp(c(-\xi_{1} + 0.5 \xi_{2} - 0.25 \xi_{3} - 0.1 \xi_{4}))}{1+\exp(c(-\xi_{1} + 0.5 \xi_{2} - 0.25 \xi_{3} - 0.1 \xi_{4}))},
\label{eq:ks_variant_c}
\end{align}
where we take $c \in \{0.25, 0.5, 0.75, 1, 1.25, 1.5, 1.75\}$. The case in which $c = 1$ matches the standard setup of \cite{KangSc07}. As discussed in \Cref{sec:ks}, the closer the value of the scaling $c$ to 0, the more concentrated the treatment probabilities are to 50\%. In the other direction, the higher the value of $c$, the more extreme propensities are observed, increasing the variance of the propensity scores, pushing the overlap assumption to the limit and making the problem of causal estimation more challenging.

For each value of the parameter $c$, we generate 1000 seeds of sample size 200 with no truncation. We use linear outcome models and logistic propensity models. The results for the mean absolute error are displayed in
Figure~\ref{fig:mae_vs_prop}. In Table~\ref{tb:prop_stats_ks} and Table~\ref{tb:1_prop_stats_ks} we show the main statistics for the fitted propensity scores $\what \pi (X)$ and $1\what \pi(X)$ as a function of the scaling parameter. The statistics are taken over all observations (treated and non-treated).
  
\begin{table}[ht]
\centering
\footnotesize
\begin{tabular}{lrrrrr}
  \toprule
Scaling $c$ & Stdev $\what\pi(X)$  & Min $\what\pi(X)$ & Max $\what\pi(X)$ & CVar $\what\pi(X)$ \\ 
  \midrule
0.25 & 0.08 (0.001) & 0.2021 (0.003) & 0.70 (0.003) & 0.296 (0.003) \\ 
  0.50 & 0.13 (0.001) & 0.0756 (0.002) & 0.77 (0.002) & 0.186 (0.002) \\ 
  0.75 & 0.17 (0.001) & 0.0240 (0.001) & 0.83 (0.002) & 0.102 (0.002)\\ 
  1.00 & 0.21 (0.001) & 0.0066 (0.000) & 0.88 (0.002) & 0.055 (0.001)\\ 
  1.25 & 0.24 (0.001) & 0.0018 (0.000) & 0.92 (0.001)& 0.029 (0.001)\\ 
  1.50 & 0.26 (0.001) & 0.0006 (0.000) & 0.94 (0.001)& 0.016 (0.001)\\ 
  1.75 & 0.28 (0.001) & 0.0002 (0.000) & 0.96 (0.001) & 0.009 (0.000) \\ 
   \bottomrule
\end{tabular}
\caption{Summary statistics of $\what\pi(X)$ for different values of scaling parameter $c$ from Equation~\eqref{eq:ks_variant_c}. We show means across 1000 dataset draws of size 200, with standard error in parentheses. CVar refers to the mean of the 5\% of smallest values.}
\label{tb:prop_stats_ks}
\end{table}
  
\begin{table}[ht]
\centering
\footnotesize
\begin{tabular}{lrrrrr}
  \toprule
Scaling $c$ & Stdev $1/\what\pi(X)$  & Min $1/\what\pi(X)$ & Max $1/\what\pi(X)$ &  CVar $1/\what\pi(X)$  \\ 
  \midrule
0.25 & 0.5 (0) & 1.42 (0.005) & 5 (1.10) & 4 (0.128)\\ 
  0.50 & 1.3 (1)& 1.30 (0.004) & 13 (7.57) & 6 (1.05) \\ 
  0.75 & 3.9 (11)& 1.20 (0.003) & 42 (148)& 15 (16.3)\\ 
  1.00 & 13.2 (1016)& 1.13 (0.002) & 152 (14325) & 41 (1453) \\ 
  1.25 & 44.0 (16634)& 1.09 (0.001) & 543 (235232) & 117  (23538)\\ 
  1.50 & 136.3 (548595) & 1.06 (0.001) & 1692 (7758318) & 348 (775940) \\ 
  1.75 & 440.1 (94410206) & 1.04 (0.001)  & 5585 (1335161973) & 1003 (133516504) \\ 
   \bottomrule
\end{tabular}
\caption{Summary statistics of $1/\what\pi(X)$ for different values of scaling parameter $c$ from Equation~\eqref{eq:ks_variant_c}. We show means across 1000 dataset draws of size 200, with standard error in parentheses. CVar refers to the mean of the 5\% of largest values.}
\label{tb:1_prop_stats_ks}
\end{table}

\subsubsection{Additional Results}
\label{sec:ks_more_results} 
\begin{table}[H]
\footnotesize
\vspace{-0.5cm}
\begin{subtable}{\linewidth}
\centering
\caption{$N=200$, no truncation} 
\begin{tabular}{lrrrrrrr}
\toprule
Method &  \multicolumn{2}{c}{Bias} &  \multicolumn{2}{c}{Mean Abs Err} &  \multicolumn{2}{c}{RMSE} &  \multicolumn{1}{c}{Median Abs Err} \\
\midrule
Direct & -0.00 & (0.103) & 2.60 & (0.061) & 3.24 & (0.457) & 2.13 \\
IPW & 22.10 & (2.579) & 27.28 & (2.52)& 84.45 & (2733) & 9.87 \\ 
IPW-SN & 3.36 &(0.292) & 5.42 & (0.260) & 9.83 & (13.45) & 3.01 \\ 
\greymidrule
AIPW & -5.08 & (0.474) & 6.16 & (0.461) & 15.82 & (87.1) & 3.43 \\ 
AIPW-SN & -3.65 & (0.203) & 4.73 & (0.179) & 7.38 & (6.14) & 3.26 \\ 
TMLE & -111.59 & (41.073) & 112.15 & (41.1) & 1302.98 & ($10^6$) & 3.95 \\ 
C-Learner & \textbf{-2.45} & (0.120) & \textbf{3.57} & (0.088) & \textbf{4.52} & (0.912) & \textbf{2.93} \\ 
\bottomrule
\end{tabular}
\end{subtable}

\begin{subtable}{\linewidth}
\centering
\caption{$N=1000$, no truncation} %
\begin{tabular}{lrrrrrrr}
\toprule
Method & \multicolumn{2}{c}{Bias} &  \multicolumn{2}{c}{Mean Abs Err} &  \multicolumn{2}{c}{RMSE} &  \multicolumn{1}{c}{Median Abs Err} \\
\midrule
Direct & -0.43 & (0.044 )& 1.17 & (0.028) & 1.46 & (0.095) & 1.00 \\
IPW & 105.46 & (59.843) & 105.67 & (59.8) & 1894.40 & ($10^6$) & 17.95 \\ 
IPW-SN & 6.83 & (0.331) & 7.02 & (0.327) & 12.50 & (24.0) & \textbf{4.09} \\
\greymidrule
AIPW & -41.37 & (24.821) & 41.39 & (24.8) & 785.60 & ($10^5$) & 5.22 \\ 
AIPW-SN & -8.35 & (0.431) & 8.37 & (0.430) & 15.97 & (46.4) & 4.92 \\
TMLE & -17.51 & (3.493) & 17.51 & (3.49) & 111.77 & ($10^4$) & 4.25 \\ 
C-Learner & \textbf{-4.40} & (0.077) & \textbf{4.42} & (0.076) & \textbf{5.03} & (0.795) & 4.21 \\ 
\bottomrule
\end{tabular}
\end{subtable}

\begin{subtable}{\linewidth}
\centering
\caption{$N=200$, truncation threshold 5\%} %
\begin{tabular}{lrrrrrrr}
\toprule
 Method & \multicolumn{2}{c}{Bias} &  \multicolumn{2}{c}{Mean Abs Err} &  \multicolumn{2}{c}{RMSE} &  \multicolumn{1}{c}{Median Abs Err} \\
\midrule
Direct & -0.00 & (0.103) & 2.60 & (0.061) & 3.24 & (0.457) & 2.13 \\ 
IPW & 5.13 & (0.398) & 10.47 & (0.275) & 13.60 & (10.055) & 8.35 \\ 
IPW-SN & \textbf{1.17} & (0.132) & 3.33 & (0.087) & 4.32 & (0.969) & \textbf{2.59} \\ 
\greymidrule
AIPW & -2.45 & (0.121) & 3.59 & (0.088) & 4.54 & (0.913) & 3.02 \\ 
AIPW-SN & -2.37 & (0.118) & 3.50 & (0.085) & 4.42 & (0.859) & 2.92 \\ 
TMLE & -2.01 & (0.107) & \textbf{3.15} & (0.075) & \textbf{3.95} & (0.665) & 2.62 \\ 
C-Learner & -2.15 & (0.112) & 3.31 & (0.079) & 4.14 & (0.734) & 2.83 \\ 
\bottomrule
\end{tabular}
\end{subtable}

\begin{subtable}{\linewidth}
\centering
\caption{$N=1000$, truncation threshold 5\%} %
\begin{tabular}{lrrrrrrr}
\toprule
 Method & \multicolumn{2}{c}{Bias} &  \multicolumn{2}{c}{Mean Abs Err} &  \multicolumn{2}{c}{RMSE} &  \multicolumn{1}{c}{Median Abs Err} \\
\midrule
Direct & -0.43 & (0.044) & 1.17 & (0.028) & 1.46 & (0.095) & 1.00 \\ 
IPW &  8.35 & (0.179) & 8.67 & (0.163) & 10.08 & (3.413) & 8.23 \\ 
IPW-SN & \textbf{1.95} & (0.062) & 2.27 & (0.050) & 2.76 & (0.293) & \textbf{2.01} \\
\greymidrule
AIPW & -3.41 & (0.053) & 3.44 & (0.051) & 3.80 & (0.378) & 3.44 \\ 
AIPW-SN &  -3.31 & (0.052) & 3.34 & (0.050) & 3.70 & (0.358) & 3.35 \\ 
TMLE & -2.81 & (0.046) & \textbf{2.85} & (0.044) & \textbf{3.17} & (0.274) & 2.86 \\ 
C-Learner & -3.07 & (0.049) & 3.10 & (0.047) & 3.43 & (0.310) & 3.10 \\ 
\bottomrule
\end{tabular}
\end{subtable}
    \caption{Estimator performance in 1000 tabular simulations for the linear specification of outcome models. ``truncation'' refers to truncation $\what\pi(X)$ away from 0. Standard error is in parentheses. Asymptotically optimal methods are listed beneath the horizontal divider. We highlight the best-performing overall method \emph{other than the direct method} in \textbf{bold}.} %
    \label{tb:linear_simulation_full}
    \vspace{-1em}
\end{table}





\begin{table}[H]
\centering
\small
\begin{tabular}{lrrrr}
\toprule
  & \multicolumn{2}{c}{$N=200$} & \multicolumn{2}{c}{$N=1000$} \\
  \midrule 
  Method & No truncation & 5\% truncation & No truncation & 5\% truncation \\ 
  \midrule
  Direct & 0.89 & 0.89 & 0.88 & 0.88 \\ 
  IPW & 1.00 & 1.00 & 1.00 & 1.00 \\ 
  IPW-SN & 1.00 & 1.00 & 1.00 & 1.00 \\ 
  \greymidrule
  AIPW & 0.88 & 0.90 & 0.59 & 0.56 \\ 
  AIPW-SN & 0.92 & 0.91 & 0.72 & 0.58 \\  
  TMLE & 0.84 & 0.90 & 0.48 & 0.60 \\ 
  C-Learner & 0.91 & 0.90 & 0.68 & 0.58 \\ 
\bottomrule
\end{tabular}
\caption{Coverage results for 1000 simulations in the linear specification. Confidence intervals were set to achieve 95\% confidence. Asymptotically optimal methods are listed beneath the horizontal divider.}
\label{tb:coverage_ks}
\end{table}

\begin{table}[H]
\footnotesize
\centering
\hfill
\begin{subtable}{.53\linewidth}
\centering
\caption{$N=200$, truncation threshold = 5\%} %
\begin{tabular}{lrrrr}
\toprule
 Method & \multicolumn{2}{c}{Bias} & \multicolumn{2}{c}{Mean Abs Err} \\
\cmidrule(lr){1-1}
\cmidrule(lr){2-3}
\cmidrule(lr){4-5}
Direct & -6.10 & (0.10) & 6.18 & (0.10)\\
IPW & 4.45 & (0.75) & 17.8 & (0.52) \\
IPW-SN & -3.40 & (0.16) & 5.06 & (0.10)\\
Lagrangian & -4.94 & (0.10) & 5.14 & (0.09) \\
\greymidrule
AIPW & \textbf{-1.80} & (0.14) & 3.85 & (0.09) \\
AIPW-SN & -2.10 & (0.12) & 3.64 & (0.08) \\
TMLE & -2.84 & (0.10) & \textbf{3.51} & (0.08) \\
C-Learner & -3.10 & (0.10) & 3.68 & (0.08) \\
\bottomrule
\end{tabular}
\end{subtable}
\hspace{-3.15em}
\begin{subtable}{.42\linewidth}
\centering
\caption{$N=200$, truncation threshold = 0.1\%} %
\begin{tabular}{rrrr}
\toprule
 \multicolumn{2}{c}{Bias} & \multicolumn{2}{c}{Mean Abs Err} \\
\cmidrule(lr){1-2}
\cmidrule(lr){3-4}
-6.10 & (0.10) & 6.18 & (0.10)\\
41 & (5.82) & 54.2 & (5.72) \\
-1.10 & (0.38) & 7.20 & (0.31)\\
-4.96 & (0.10) & 5.16 & (0.09) \\
\greymidrule
4.82 & (1.27) & 10.4 & (1.23) \\
\textbf{-0.57} & (0.26) & 5.02 & (0.21) \\
-1.42 & (0.17) & 4.19 & (0.12) \\
-3.24 & (0.09) & \textbf{3.79} & (0.08) \\
\bottomrule
\end{tabular}
\end{subtable}
\hfill
\hspace*{\fill}
\caption{Comparison of estimator performance on misspecified datasets from \citet{KangSc07} in 1000 tabular simulations using gradient boosted regression trees with 200 samples, 5\%, and 0.1\% truncation threshold. Asymptotically optimal methods are listed beneath the horizontal divider. We highlight the best-performing method in \textbf{bold}. Standard errors are displayed within parentheses to the right of the point estimate. ``Lagrangian'' refers to only performing the first stage in Algorithm~\ref{alg:boosting}.
}

    \label{tb:linear_simulation_xgb_trim}
\end{table}





\begin{table}[H]
\footnotesize
\hfill
\hspace*{\fill}
\begin{subtable}{.42\linewidth}
\centering
\caption{$N=200$}
\begin{tabular}{lrrrr}
\toprule
Method & \multicolumn{2}{c}{Bias} & \multicolumn{2}{c}{Mean Abs Err} \\
\cmidrule(lr){1-1}
\cmidrule(lr){2-3}
\cmidrule(lr){4-5}
Direct & 4.82 & (0.091) & 4.94 & (0.085) \\ 
IPW & \textbf{-0.70} & (0.141) & 3.44 & (0.092)  \\ 
IPW-SN & 2.48 & (0.097) & \textbf{3.23} & (0.072)  \\
    \greymidrule
AIPW & 3.23 & (0.093) & 3.63 & (0.077)  \\ 
AIPW-SN & 3.22 & (0.092) & 3.61 & (0.076)  \\ 
TMLE & 3.11 & (0.091) & 3.50 & (0.076)  \\ 
C-Learner & 3.19 & (0.092) & 3.58 & (0.076)  \\ 
\bottomrule
\end{tabular}
\end{subtable}
\begin{subtable}{.42\linewidth}
\centering
\caption{$N=1000$} %
\begin{tabular}{rrrr}
\toprule
 \multicolumn{2}{c}{Bias} & \multicolumn{2}{c}{Mean Abs Err} \\
\cmidrule(lr){1-2}
\cmidrule(lr){3-4}
 4.80 & (0.041) & 4.80 & (0.041)  \\ 
 \textbf{-0.83} & (0.054) & \textbf{1.52} & (0.036)  \\  
2.43 & (0.042) & 2.47 & (0.040)  \\
    \greymidrule
3.14 & (0.041) & 3.15 & (0.040) \\ 
 3.12 & (0.041) & 3.13 & (0.040)  \\ 
 3.03 & (0.041) & 3.04 & (0.040)  \\ 
3.09 & (0.041) & 3.10 & (0.040) \\
\bottomrule
\end{tabular}
\end{subtable}
\hfill
\hspace*{\fill}
\caption{Results of 1000 simulations for estimating $P[Y(0)]$ instead of $P[Y(1)]$  with the linear specification. The standard error is in parentheses. Asymptotically optimal methods are listed beneath the horizontal divider.}
\label{tb:flipped_ks}
\end{table}








\begin{table}[H]
\centering
\small
\begin{tabular}{lrrrr}
\toprule
 \multicolumn{1}{c}{Method} & \multicolumn{2}{c}{$N=200$} & \multicolumn{2}{c}{$N=1000$} \\
 \cmidrule(lr){1-1}
\cmidrule(lr){2-3}
\cmidrule(lr){4-5}
  Direct & 0.35 & (0.01) & 0.12 & (0.01) \\ 
  IPW & 1.00 & (0.00) & 0.98 & (0.01) \\ 
  IPW-SN & 0.85 & (0.01) & 1.00 & (0.00) \\ 
Lagrangian & 0.84 & (0.01) & 0.71 & (0.01) \\ 
  \greymidrule
  AIPW & 0.85 & (0.01) & 0.85 & (0.01) \\ 
  AIPW-SN & 0.85 & (0.01) & 0.85 & (0.01) \\  
  TMLE & 0.86 & (0.01) & 0.87 & (0.01) \\ 
  C-Learner & 0.82 & (0.01) & 0.77 & (0.01) \\ 
\bottomrule
\end{tabular}
\caption{Coverage results for 1000 simulations using gradient boosted regression trees. Confidence intervals were set to achieve 95\% confidence. Asymptotically optimal methods are listed beneath the horizontal divider. ``Lagrangian'' refers to only performing the first stage in Algorithm~\ref{alg:boosting}.} %
\label{tb:coverage_ks_xgb}
\end{table}




















\paragraph{Balancing results, linear outcome model}
\ifdefined\response\newpage\newpage\clearpage\pagebreak\else\fi
\label{sec:linear_cbps_ks}
\begin{table}[t]
\footnotesize
\centering
\hfill
\begin{subtable}{.53\linewidth}
\centering
\caption{$N=200$} %
\begin{tabular}{lrrrr}
\toprule
 Method & \multicolumn{2}{c}{Bias} & \multicolumn{2}{c}{Mean Abs Err} \\
\cmidrule(lr){1-1}
\cmidrule(lr){2-3}
\cmidrule(lr){4-5}
CBPS-IPW & -6.92 & (0.26) & 8.52 & (0.21) \\
CBPS-IPW-SN & -2.98 & (0.10) & 3.66 & (0.08) \\
\greymidrule
CBPS-AIPW & -1.68 & (0.36) & 3.11 & (0.24) \\
CBPS-AIPW-SN &  -1.69 & (0.36) & 3.11 & (0.24) \\ 
CBPS-TMLE &  -7.20 & (5.14) & 9.50 & (5.10) \\ 
CBPS-C-Learner &  -1.70 & (0.36) & 3.10 & (0.24) \\ 
\greymidrule
CBPS-TMLE-L & -2.01 & (0.33) & 3.07 & (0.24) \\ 
CBPS-C-Learner-L & \textbf{-1.85} & (0.34) & \textbf{2.97} & (0.24) \\ 
\bottomrule
\end{tabular}
\end{subtable}
\hspace{-3.15em}
\begin{subtable}{.42\linewidth}
\centering
\caption{$N=1000$} %
\begin{tabular}{rrrr}
\toprule
 \multicolumn{2}{c}{Bias} & \multicolumn{2}{c}{Mean Abs Err} \\
\cmidrule(lr){1-2}
\cmidrule(lr){3-4}
1.76 & (0.15) & 4.00 & (0.10) \\
\textbf{-1.55} & (0.06) & \textbf{1.88} & (0.04) \\
\greymidrule
-3.54 & (0.37) & 3.58 & (0.37) \\ 
-3.44 & (0.33) & 3.49 & (0.33) \\ 
-7.54 & (2.48) & 7.55 & (2.48) \\
-3.10 & (0.20) & 3.15 & (0.19) \\ 
\greymidrule
-2.76 & (0.15) & 2.79 & (0.14) \\ 
-2.76 & (0.18) & 2.83 & (0.17) \\ 
\bottomrule
\end{tabular}
\end{subtable}
\hfill
\hspace*{\fill}
\caption{%
Comparison of estimator performance on misspecified datasets from \citet{KangSc07} in 1000 tabular simulations, for linear outcome models (\Cref{sec:ks_linear}) using propensity scores fitted via covariate balancing. We include only methods that makes use of propensity scores.  We highlight the best-performing method in \textbf{bold}. Standard errors are displayed within parentheses to the right of the point estimate.
}
    \label{tb:linear_cbps}
\end{table}

As discussed in \Cref{sec:background}, covariate balancing techniques imply constraints, possibly via customized loss functions, to equate the covariates of samples in treatment and control for a given propensity score model $\hat \pi$, thus, informing how to constraint or estimate the propensity model $\hat \pi$. Naturally, methods that make use of propensity score such as IPW (the basic CBPS approach), AIPW, TMLE and C-Learner, may also benefit from a tailored propensity score function $\hat \pi$. 

Therefore, we also provide in \Cref{tb:linear_cbps} a comparison with the Covariate Balancing Propensity Score (CBPS). We discuss the implementation details of CBPS in \Cref{sec:ks_experiment_details}. Roughly speaking, all asymptotically optimal methods improve or maintain its performance when using propensity scores via covariate balancing, despite an increase in the MAE standard error. For the sample size of 200, C-Learner with covariate balancing and logistic loss is the best performance method for the MAE point estimate, although most of the asymptotically optimal methods are statistically the same. When the sample size is 1000, covariate balancing with self-normalization is the best performing method. %
\label{sec:linear_cbps_ks-end}
\ifdefined\response\newpage\newpage\clearpage\pagebreak\else\fi




\subsection{CivilComments Experiments (Section~\ref{sec:civilcomments})}
\subsubsection{Implementation Details}
\label{sec:nlp_details}


\paragraph{Model and Training Objectives}
Both propensity and outcome models are a linear layer on top of a DistilBERT featurizer, with either softmax and cross-entropy loss for propensity score, or mean squared error for outcome models. 
For training outcome models, we only consider training data on which $A=1$, as only those terms contribute to the loss. 
\ifdefined\pnas
    $\what\mu^C$ learned using the procedure described in \cref{sec:methods_nn}. 
    We train $\what \mu$
    using squared loss on labeled units, and propensity models $\what \pi$ using the logistic loss. 
    On each dataset draw, we use cross-fitting, as described in \cref{sec:theory_text}, with $K=2$ folds, for all estimators. 
\fi
\paragraph{Hyperparameters and Settings} 
For propensity models, we tried learning rates of $\{10^{-3},10^{-4},10^{-5}\}$
for the setting of $l=10^{-4}$ and found that a learning rate of $10^{-4}$ performed the best in terms of val loss. Because of computational constraints, we also used this learning rate for $l=10^{-2},10^{-3}$. 

For outcome models (including C-Learner) in the settings of $l\in\{10^{-2},10^{-3},10^{-4}\}$, we tried learning rates of $\{10^{-3},10^{-4},10^{-5},10^{-6}\}$. 
For C-Learner regularization, we tried $\lambda$ taking on values of $\lambda_0 / P_{\rm eval}[A/\hat\pi(X)]^2$ where $\lambda_0\in\{0,1,4,16,64,256\}$. 
Essentially, the penalty is $\lambda_0$ times the square of the bias shift that would be required to satisfy the condition. 

Because of computational constraints, for choosing these hyperparameters (learning rate, $\lambda$), we select a subset of dataset draws, and choose the best hyperparameters based on each model's criteria on that small set of dataset draws. Then, after the hyperparameters are chosen, we run over the entire set of dataset draws. 

While hyperparameters for the (usual) outcome model are chosen to minimize val loss, in contrast, the hyperparameters (learning rate, regularization $\lambda$) for the C-Learner outcome model were selected to minimize the \emph{magnitude of the constant shift} at the end of the first epoch, as in \cref{sec:methods_nn}. 
The idea is that if the size of the constant shift is small, then the regularizer is doing a good job of enforcing the constraint. 
Although one might think that larger regularization $\lambda$'s would automatically lead to smaller constant shifts, we do not find this to be the case; the best $\lambda$ in our settings (64) is not the largest value that we try (up to 256). 
Furthermore, the set of hyperparameter values we choose between in this process are ones that seem to result in reasonable performance, e.g. in terms of MSE on treated units in the validation set. 


The best hyperparameters chosen are in Table~\ref{tab:cc_hparams}. 
We trained all models with batch size 64. Learning rates decayed linearly over 10 epochs. Optimization was done using the AdamW optimizer. Weight decay was fixed at $0.1$. 

We choose the training epoch with the best val loss for each model (cross entropy for propensity, and MSE for outcome). For C-Learner, we choose the epoch with the best val MSE. 



\subsubsection{Additional Results}
\label{sec:more_nlp_results}
We consider three data generating processes as in \cref{sec:civilcomments}, where treatment assignment is parameterized by $l=10^{-4},10^{-3},10^{-2}$, respectively. 
Summary statistics of $\what\pi(X)$ are in \Cref{tab:civilcomments_prop}, and summary statistics of $1/\what\pi(X)$ are in \cref{tab:civilcomments_invprop}.
Hyperparameters for outcome models are in Table~\ref{tab:cc_hparams}. 
Comparisons of estimators in these settings are in Tables~\ref{tab:nn_4},~\ref{tab:nn_3},~\ref{tab:nn_2}.

\begin{table}[H]
\centering
\footnotesize
\begin{tabular}{lrrrr}
\toprule
Value of $l$ & Min $\what \pi(X)$ & Max $\what \pi(X)$ & Std $\what \pi(X)$ & CVar $\what \pi(X)$ \\
\midrule 
$l=10^{-2}$ &0.0014 (0.0001) &         0.94 (0.01) &         0.15 (0.002) &         0.0016 (0.0001)         \\
$l=10^{-3}$ &0.0004 (0.0000) &         0.95 (0.01) &         0.15 (0.002) &         0.0004 (0.0000)         \\
$l=10^{-4}$ &0.0003 (0.0000) &         0.96 (0.00) &         0.16 (0.002) &         0.0003 (0.0000)         \\
\bottomrule
\end{tabular}
\caption{Summary statistics of $\what\pi(X)$ for values of hyperparameter $l=10^{-4},10^{-3},10^{-2}$ for \cref{sec:civilcomments} experiments. We show means across 100 dataset draws for each value of $l$, with standard error in parentheses. CVar refers to the mean of the 5\% of smallest values. }
    \label{tab:civilcomments_prop}
\end{table}
\begin{table}[H]
\centering
\footnotesize
\begin{tabular}{lrrrr}
\toprule
Value of $l$ & Min $1/\what \pi(X)$ & Max $1/\what \pi(X)$ & Std $1/\what \pi(X)$ & CVar $1/\what \pi(X)$ \\
\midrule 
$l=10^{-2}$ &1.07 (0.01) &         1124.7\;\; (91.0) &         237.1 (23.4) &         1050.1 
 \; (85.3) \\
$l=10^{-3}$ &1.06 (0.01) &         3244.0 (190.5) &         757.7 (54.1) &         3090.9 (183.5) \\
$l=10^{-4}$ &1.05 (0.00) &         4126.9 (214.6) &         980.6 (61.4) &         3960.2 (205.7) \\
\bottomrule
\end{tabular}
\caption{Summary statistics of $1/\what\pi(X)$ for values of hyperparameter $l=10^{-4},10^{-3},10^{-2}$ for \cref{sec:civilcomments} experiments. We show means across 100 dataset draws for each value of $l$, with standard error in parentheses. CVar refers to the mean of the 5\% of largest values. }
    \label{tab:civilcomments_invprop}
\end{table}
\begin{table}[H]
\centering
\footnotesize
\begin{tabular}{lrrrrrr}
\toprule
 & $l=10^{-4}$ & $l=10^{-3}$ & $l=10^{-2}$ \\
 \midrule 
Outcome model learning rate & $10^{-3}$ & $10^{-3}$ & $10^{-4}$ \\
C-Learner learning rate (best val MSE) & $10^{-4}$ & $10^{-4}$ & $10^{-4}$ \\ 
C-Learner learning rate (min bias shift) & $10^{-5}$ & $10^{-5}$ & $10^{-5}$ \\ 
C-Learner $\lambda$ (best val MSE) & 0 & 16 & 16 \\
C-Learner $\lambda$ (min bias shift) & 64 & 64 & 64 \\
\bottomrule
\vspace{0.05em}
\end{tabular}
\caption{Hyperparameters for $l=10^{-4},10^{-3},10^{-2}$ for \cref{sec:civilcomments}}
    \label{tab:cc_hparams}
\end{table}

\begin{table}[H]
\centering
\footnotesize
\begin{tabular}{lrrrrrr}
\toprule
Method & \multicolumn{2}{c}{Bias} & \multicolumn{2}{c}{Mean Abs Err} & \multicolumn{2}{c}{Coverage}\\ 
\cmidrule(lr){1-1}
\cmidrule(lr){2-3}
\cmidrule(lr){4-5}
\cmidrule(lr){6-7}
Direct & 0.173 & (0.008) & 0.177 & (0.007) & 0.010 & (0.001)\\
IPW & 0.504 & (0.084) & 0.546 & (0.081) & 0.760 & (0.018)\\
IPW-SN & 0.114 & (0.017) & 0.153 & (0.014) & 0.890 & (0.010)\\
\greymidrule
AIPW & 0.084 & (0.043) & 0.307 & (0.032) & 0.830 & (0.014)\\
AIPW-SN & 0.116 & (0.018) & 0.161 & (0.014) & 0.850 & (0.013)\\
TMLE & -1.264 & (1.361) & 1.802 & (1.355) & 0.810 & (0.015)\\
C-Learner (best val MSE) & 0.103 & (0.015) & 0.141 & (0.011) & 0.870 & (0.011)\\
C-Learner (min bias shift) & \textbf{0.075} & (0.012) & \textbf{0.115} & (0.008) & 0.900 & (0.009)\\
\bottomrule
\vspace{0.05em}
\end{tabular}
\caption{Comparison of estimators in the CivilComments~\citep{civilcomments} semi-synthetic dataset over 100 re-drawn datasets, with $l=10^{-4}$. Confidence intervals were set to achieve 95\% confidence. Asymptotically optimal methods are listed beneath the horizontal divider. We highlight the best-performing overall method in \textbf{bold}. Standard errors are displayed within parentheses to the right of the point estimate.}
    \label{tab:nn_4}
\end{table}


\begin{table}[H]
\centering

\footnotesize
\begin{tabular}{lrrrrrr}
\toprule
Method & \multicolumn{2}{c}{Bias} & \multicolumn{2}{c}{Mean Abs Err} & \multicolumn{2}{c}{Coverage}\\ 
\cmidrule(lr){1-1}
\cmidrule(lr){2-3}
\cmidrule(lr){4-5}
\cmidrule(lr){6-7}
Direct & 0.147 & (0.028) & 0.200 & (0.025) & 0.000 & (0.000)\\
IPW & 0.417 & (0.064) & 0.437 & (0.063) & 0.840 & (0.013)\\
IPW-SN & 0.079 & (0.015) & 0.120 & (0.013) & 0.910 & (0.008)\\
\greymidrule
AIPW & \textbf{-0.056} & (0.052) & 0.344 & (0.039) & 0.950 & (0.005)\\
AIPW-SN & 0.089 & (0.020) & 0.134 & (0.018) & 0.950 & (0.005)\\
TMLE & 0.074 & (0.025) & 0.144 & (0.022) & 0.940 & (0.006)\\
C-Learner (best val MSE) & 0.067 & (0.012) & 0.110 & (0.009) & 0.980 & (0.002)\\
C-Learner (min bias shift) & 0.060 & (0.011) & \textbf{0.099} & (0.008) & 0.990 & (0.001)\\
\bottomrule
\vspace{0.05em}
\end{tabular}
\caption{Comparison of estimators in the CivilComments~\citep{civilcomments} semi-synthetic dataset over 100 re-drawn datasets, with $l=10^{-3}$. Confidence intervals were set to achieve 95\% confidence. Asymptotically optimal methods are listed beneath the horizontal divider. We highlight the best-performing overall method in \textbf{bold}. Standard errors are displayed within parentheses to the right of the point estimate.}
    \label{tab:nn_3}
\end{table}



\begin{table}[H]
\centering

\footnotesize
\begin{tabular}{lrrrrrr}
\toprule
Method & \multicolumn{2}{c}{Bias} & \multicolumn{2}{c}{Mean Abs Err} & \multicolumn{2}{c}{Coverage}\\ 
\cmidrule(lr){1-1}
\cmidrule(lr){2-3}
\cmidrule(lr){4-5}
\cmidrule(lr){6-7}
Direct & 0.091 & (0.007) & 0.099 & (0.006) & 0.040 & (0.004)\\
IPW & 0.346 & (0.041) & 0.353 & (0.040) & 0.710 & (0.021)\\
IPW-SN & 0.004 & (0.005) & 0.038 & (0.003) & 0.950 & (0.005)\\
\greymidrule
AIPW & -0.172 & (0.029) & 0.220 & (0.026) & 0.890 & (0.010)\\
AIPW-SN & 0.002 & (0.005) & \textbf{0.040} & (0.003) & 1.000 & (0.000)\\
TMLE & 0.027 & (0.005) & 0.042 & (0.003) & 0.990 & (0.001)\\
C-Learner (best val MSE) & 0.009 & (0.007) & 0.053 & (0.004) & 1.000 & (0.000)\\
C-Learner (min bias shift) & \textbf{0.001} & (0.006) & 0.045 & (0.003) & 1.000 & (0.000)\\
\bottomrule
\vspace{0.05em}
\end{tabular}
\caption{Comparison of estimators in the CivilComments~\citep{civilcomments} semi-synthetic dataset over 100 re-drawn datasets, with $l=10^{-2}$. Confidence intervals were set to achieve 95\% confidence. Asymptotically optimal methods are listed beneath the horizontal divider. We highlight the best-performing \emph{asymptotically optimal} method in \textbf{bold}. Standard errors are displayed within parentheses to the right of the point estimate. C-Learner's performance is within standard error of AIPW-SN and TMLE.}
    \label{tab:nn_2}
\end{table}
C-Learner errors are comparatively more stable for $l=10^{-4},10^{-3}$. For $l=10^{-2}$, C-Learner is comparable to other asymptotically optimal methods. 




\ifdefined\msom\else
\subsection{IHDP Tabular Dataset}
\label{sec:ihdp_experiment}
The Infant Health and Development Program (IHDP) dataset is a tabular semisynthetic dataset~\cite{brooks1992effects} that was first introduced as a benchmark for ATE estimation by \citet{hill2011bayesian}. IHDP is based on a randomized experiment studying the effect of specialized healthcare interventions on the cognitive scores of premature infants with low birth weight. The dataset consists of a continuous outcome variable $Y$, a binary treatment variable $A$, and 25 covariates $X$ that affect the outcome variable and are also correlated with the treatment assignment. We use 1000 datasets %
as in %
\cite{shi2019adapting,chernozhukov2022riesznet}.

\subsubsection{Gradient Boosted Regression Trees}
\label{sec:ihdp_xgb}
We instantiate C-Learner using gradient boosted regression trees using the XGBoost package \citep{xgb} with a custom objective, as outlined in \cref{sec:methods_boosting}. $\what\pi$ is fit as a logistic regression on covariates $X$. %
Different to the applications presented in the main text, we are interested in estimating the ATE $P[Y(1) - Y(0)]$ where $Y(0)$ is nonzero. We discuss two alternatives to instantiate the C-Learner. First, we find an outcome model that takes as input both covariates and the treatment variable, which leads to the following formulation
\begin{equation*}\label{eq:c-learner-full}
\what \mu^C_{-k,n}=\argmin_{\tilde\mu\in\mathcal F} 
\left\{ P_{-k,n}[(Y-\tilde\mu(A,X))^2] : 
P_{k,n}\left[\left(\frac{A}{\what \pi(X)}-\frac{1-A}{1-\what \pi(X)}\right)(Y-\tilde\mu(A,X))\right]=0
\right\},
\end{equation*}
and the final estimator becomes
\begin{align*}
    \what\psi^C_n :=\frac{1}{K}\sum_{k=1}^K P_{k,n}[\what\mu_{-k,n}^C(1,X)-\what\mu_{-k,n}^C(0,X)].
\label{eq:c_learner_def_xgb}
\end{align*}
This approach is often referred to as the ``S-Learner".


Another alternative is to model the ATE as the difference of two missing outcomes and fit two outcome models for treated vs. non-treated units. This approach is commonly referred to as the ``T-Learner". In this case, we estimate two models by solving
\begin{align*}
\what{\mu}_{-k, n, 1}^C 
& \in \argmin_{\tilde{\mu} \in \mc{F}}
\left\{ P_{-k, n} [ A (Y - \tilde{\mu}(X))^2]:
P_{k,n}\left[\frac{A}{\what\pi_{-k,n}(X)}(Y-\tilde{\mu}(X))\right]=0
\right\}, \\
\what{\mu}_{-k, n, 0}^C 
& \in \argmin_{\tilde{\mu} \in \mc{F}}
\left\{ P_{-k, n} [ (1-A) (Y - \tilde{\mu}(X))^2]:
P_{k,n}\left[\frac{1-A}{1-\what\pi_{-k,n}(X)}(Y-\tilde{\mu}(X))\right]=0
\right\},
\end{align*}
and the final estimator becomes
\begin{align*}
    \what\psi^C_n :=P_{\rm eval}\frac{1}{K}\sum_{k=1}^K P_{k,n}[\what\mu_{-k,n,1}^C(X)-\what\mu_{-k,n,0}^C(X)].
\end{align*}
We empirically observe that the latter approach (T-Learner) is critical for performance in this setting: using a simple XGBoost T-Learner model significantly improves upon sophisticated S-Learner variants, including state-of-the-art methods such as RieszNet, RieszForest, and DragonNet~\cite{chernozhukov2022riesznet,shi2019adapting}.
 
For consistency with prior work ~\cite{van2006targeted,chernozhukov2022riesznet,shi2019adapting},
for each dataset run, we split the data in 80\% training and 20\% for hyperparameter tuning, and use the whole dataset when evaluating the first-order error term, i.e.,  $P_{\rm eval} = P_{\rm train} \cup P_{\rm val}$. For each outcome model, $\what{\mu}_{-k, n, 1}^C$ and $\what{\mu}_{-k, n, 0}^C$, we follow the same setup described in the Kang \& Schafer experiment. Hyperparameter tuning is performed using a grid search for the following parameters: learning rate (0.01, 0.05, 0.1, 0.2), feature subsample by tree (0.5, 0.8, 1), and max tree depth (3, 4, 5). We also set the maximum number of weak learners to 2000, and we perform early stopping using MSE loss on $P_{\rm val}$ for 20 rounds. Hyperparameter tuning is performed separately for C-Learner and the initial outcome model. For the second phase of Algorithm ~\ref{alg:boosting}, we use the set of hyperparameters found in the first stage. The weak learners are fitted using $P_{\rm eval}$. In order to avoid overfitting in the targeting step, we use a subsampling of 50\% and early stopping after 20 rounds.



 The results for the mean absolute error and their respective standard errors are displayed in Table~\ref{tb:ihdp}. %
The C-Learner improves upon the direct method while using an identical model class.
The asymptotically optimal estimators perform similarly, and C-Learner's advantage over other asymptotically optimal methods is not statistically significant. This is perhaps not surprising in this setting as the propensity weights do not vary as much here, in contrast to Kang \& Shafer's example where existing asymptotically optimal methods performed poorly as estimated propensity weights varied dramatically. 

\begin{table}[t]
\centering

\footnotesize
\begin{tabular}{lrrrr}
\toprule
Method & \multicolumn{2}{c}{Bias} & \multicolumn{2}{c}{Mean Abs Err} \\ 
\cmidrule(lr){1-1}
\cmidrule(lr){2-3}
\cmidrule(lr){4-5}
Direct (Boosting) & -0.16 & (0.004) & 0.188 & (0.004) \\
IPW & -1.50 & (0.040) & 1.500 & (0.040) \\ 
IPW-SN & -0.01 & (0.003) & 0.110 & (0.003) \\ 
Lagrangian & -0.09 & (0.004) & 0.144 & (0.004) \\
\greymidrule
AIPW & -0.05 & (0.002) & 0.103 & (0.003) \\ 
AIPW-SN & -0.04 & (0.002) & \textbf{0.103} & (0.003) \\ 
TMLE  & 0.007 & (0.003) & 0.103 & (0.003) \\ 
C-Learner & \textbf{-0.004} & (0.003) & 0.104 & (0.003) \\
\bottomrule
\vspace{0.1em}
\end{tabular}
\caption{Comparison of gradient boosted regression tree-based estimators in the IHDP semi-synthetic dataset over 1000 simulations. Asymptotically optimal methods are listed beneath the horizontal divider. We highlight the best-performing method in \textbf{bold}. Standard errors are displayed within parentheses to the right of the point estimate. ``Lagrangian'' refers to only performing the first stage in Algorithm~\ref{alg:boosting}. C-Learner’s performance is
within standard error of AIPW-SN.}
\label{tb:ihdp}
\end{table}






\subsubsection{Neural Networks}
\label{sec:ihdp_nn}
We also instantiate C-Learner as neural networks (\cref{sec:methods_nn}) on the IHDP dataset. 
In particular, we demonstrate how C-Learner can use Riesz representers (\ref{sec:other_estimands}), or equivalently, propensity models in the ATE setting, that are
learned by other methods.
This demonstrates the versatility of C-Learner as it is able to leverage new methods for learning Riesz representers. 
In Table~\ref{tb:ihdp_nn}, we demonstrate C-Learner using Riesz representers
learned by RieszNet~\citep{chernozhukov2022riesznet}. C-Learner achieves very similar performance to RieszNet while using the same Riesz representer. All methods use an outcome model trained in the usual way, except for RieszNet and C-Learner. In all settings where applicable, we use the Riesz representer ($A/\pi(X)-(1-A)/(1-\pi(X))$ learned by RieszNet.

We emphasize that this is not meant to be a performance comparison between RieszNet and C-Learner, for two reasons: first, we present C-Learner as a debiasing framework that is compatible with other methods for causal inference, including RieszNet; and second, we found that propensity scores fitted to the IHDP dataset are not extreme, so we expect C-Learner to perform as well as existing methods, but not to outperform existing methods. 

\begin{table}[H]
\centering
\footnotesize
\begin{tabular}{lrrrrrr}
\toprule
Method & \multicolumn{2}{c}{Bias} & \multicolumn{2}{c}{Mean Abs Err} & \multicolumn{2}{c}{Coverage} \\ 
\cmidrule(lr){1-1}
\cmidrule(lr){2-3}
\cmidrule(lr){4-5}
\cmidrule(lr){6-7}
Direct & -0.005 & (-0.000) & 0.118 & (0.003) & 0.783 & (0.025)\\
IPW & -0.789 & (-0.025) & 0.903 & (0.034) & 0.455 & (0.014)\\
IPW-SN & -0.449 & (-0.014) & 0.654 & (0.035) & 0.819 & (0.026)\\
\greymidrule
AIPW & -0.044 & (-0.001) & 0.106 & (0.003) & 0.940 & (0.030)\\
AIPW-SN & -0.042 & (-0.001) & 0.106 & (0.003) & 0.961 & (0.030)\\
RieszNet (``DR'') & \textbf{-0.033} & (-0.001) & \textbf{0.098} & (0.003) & 0.972 & (0.031)\\
C-Learner & -0.038 & (-0.001) & 0.098 & (0.002) & 0.955 & (0.030)\\
\bottomrule
\vspace{0.1em}
\end{tabular}
\caption{Comparison of neural network based estimators in the IHDP semi-synthetic dataset over 1000 simulations. Asymptotically optimal methods are listed beneath the horizontal divider. We highlight the best-performing \emph{asymptotically optimal} method in \textbf{bold}. Standard errors are displayed within parentheses to the right of the point estimate. All methods use an outcome model trained in the usual way, except for RieszNet and C-Learner. In all settings where applicable, we use the Riesz representer learned by RieszNet.}
\label{tb:ihdp_nn}
\end{table}

\paragraph{Training Procedure}
We use lightly modified RieszNet code \citep{chernozhukov2022riesznet} to generate the Riesz representers used in all methods. 
We use the same neural network architecture for RieszNet outcome models for C-Learner. 
We use $P_{\rm train}, P_{\rm val}, P_{\rm eval}$ as in \cref{sec:ihdp_xgb}. 
Hyperparameters (learning rate, $\lambda$) are selected for best mean squared error on $P_{\rm val}$ for each individual dataset. 
For the training objective, we largely follow \cref{sec:methods_nn}, with the following modifications for the full ATE rather than assuming that $Y(0)=0$: we use the mean squared error across the full dataset (as it was originally for RieszNet), and we have separate penalties and bias shifts for $A=1$ and $A=0$. 
As in RieszNet training, there is a phase of ``pre-training'' with a higher learning rate, and then a phase with a lower learning rate; for ``pre-training'' we sweep over learning rates of $\{10^{-3},10^{-4}\}$ for 100 epochs; regular training uses a learning rate of $10^{-5}$ for 600 epochs. Both use the Adam optimizer. For $\lambda$ we sweep over values of 
$\{0,0.01,0.02,0.04,0.1,0.2,0.4,1,2,4,8\}$.
Our implementation of RieszNet as a baseline is almost identical to that of the original paper, except for how we set the random seed, for initializing neural networks and for choosing train and test split. 
    

\paragraph{Comparisons with results in RieszNet paper} 
Surprisingly, our mean absolute error results in \cref{sec:ihdp_nn},
including the ones using regular RieszNet (including their algorithm code and hyperparameters), often outperform the results reported in the original RieszNet paper \citep{chernozhukov2022riesznet}. 
One difference between our code and the RieszNet code is how random seeds were set. The random seed affects the train/test split, and also weight initialization for the neural network. We re-ran the original RieszNet code, and then again with a fixed seed (set to 123) for every dataset, and we found that the results were noticeably different. 
In particular, the RieszNet mean absolute error is smaller with seeds set the second way. 
Compare Table~\ref{tab:ihdp_orig} and Table~\ref{tab:ihdp_seeds}. 
In our experiments on this setting, we fix seeds in the second way. 
\begin{table}[H]
\centering
\footnotesize
\begin{tabular}{lr}
\toprule
Method & Mean absolute error (s.e.) \\
 \midrule 
Doubly robust & 0.115 (0.003) \\
Direct & 0.124 (0.004) \\
IPW & 0.801 (0.040) \\
\bottomrule
\vspace{0.05em}
\end{tabular}
\caption{Original RieszNet IHDP results}
    \label{tab:ihdp_orig}
\end{table}

\begin{table}[H]
\centering
\footnotesize
\begin{tabular}{lr}
\toprule
Method & Mean absolute error (s.e.) \\
 \midrule 
Doubly robust & 0.098 (0.003) \\
Direct & 0.118 (0.003) \\
IPW & 0.903 (0.034) \\
\bottomrule
\vspace{0.05em}
\end{tabular}
\caption{RieszNet IHDP results, with fixed seeds}
    \label{tab:ihdp_seeds}
\end{table}


\paragraph{Relative Performance}
We've seen in \cref{sec:civilcomments} that C-Learner tends to outperform one-step estimation and targeting in settings with low overlap. Here, it merely performs about the same as one-step estimation and targeting. It seems as though overlap is not very low in IHDP, however, based on propensity scores from trained models not taking on extreme values. 
Using RieszNet's learned Riesz representers, the smallest value for $\textrm{min}(\what\pi,1-\what\pi)$ across all IHDP datasets is 0.11, which is a fair bit of overlap, especially in comparison to, say, the settings with low overlap in \cref{sec:civilcomments} where C-Learner outperformed one-step estimation and targeting methods. 
Thus, here it is expected that C-Learner matches the performance of other methods. We further emphasize that C-Learner is compatible with methods such as RieszNet, making C-Learner generally applicable.
To recap, over all of our experiments, C-Learner appears to perform either comparably to or better than one-step estimation and targeting, and outperform one-step estimation and targeting in settings with less overlap. 




\fi
